\title{DeepSeekMath: Pushing the Limits of Mathematical Reasoning in Open Language Models}

\begin{abstract}
Mathematical reasoning remains a substantial hurdle for language models because of its inherent complexity and structured demands. In this work, we present DeepSeekMath~7B, a model obtained by continuing the pre-training of DeepSeek-Coder-Base-v1.5 7B on an additional 120B tokens focused on mathematics, drawing from Common Crawl as well as incorporating both natural language and code resources. DeepSeekMath~7B attains a notable 51.7\% on the highly challenging MATH benchmark at the competition level, accomplishing this without the aid of any external toolkits or ensemble voting approaches, and drawing near the scores of models such as Gemini-Ultra and GPT-4. Furthermore, by employing self-consistency over 64 samples, DeepSeekMath~7B achieves a score of 60.9\% on the MATH dataset. The enhanced mathematical reasoning prowess of DeepSeekMath~is ascribed to two main innovations: firstly, the systematic exploitation of high-quality web-sourced data via a rigorously curated data selection workflow; and secondly, the development of Group Relative Policy Optimization (GRPO)—an adaptation of Proximal Policy Optimization (PPO)—which bolsters mathematical reasoning skill while simultaneously improving the memory efficiency compared to traditional PPO implementations.
\end{abstract}

\section{Introduction}

Large language models (LLMs) have fundamentally transformed how artificial intelligence tackles mathematical reasoning, leading to remarkable progress on quantitative reasoning benchmarks \citep{MATH} as well as geometry reasoning benchmarks \citep{trinh2024solving}. In addition to their performance gains, LLMs have also shown utility in supporting humans with challenging mathematical problem-solving tasks \citep{tao}. Despite these achievements, state-of-the-art models such as GPT-4 \citep{gpt4} and Gemini-Ultra \citep{gemini} remain inaccessible to the public, and the best available open-source alternatives exhibit a substantial performance gap.

This work presents DeepSeekMath, a specialized language model that demonstrates marked advancement in mathematical reasoning compared to existing open-source systems and attains results comparable to GPT-4 on standard academic tasks. Our approach centers around the creation of the DeepSeekMath Corpus—a large-scale, meticulously curated pre-training dataset consisting of 120B mathematical tokens. The corpus is harvested from Common Crawl (CC) via a classifier built on fastText \citep{joulin2016fasttext}. For the initial round, we train this classifier using data from OpenWebMath \citep{openwebmath} as positive samples, complemented by a representative variety of unrelated web documents serving as negative cases. With the classifier thus established, we identify further positive samples from CC, which are subsequently validated and refined with the help of human annotators. The classifier is iteratively retrained on this enriched set to boost its discriminative accuracy. Empirical assessments underscore the high caliber of the assembled corpus: our foundational DeepSeekMath-Base 7B achieves 64.2\% on GSM8K \citep{gsm8k} and 36.2\% on the MATH competition dataset \citep{MATH}, surpassing the performance of Minerva 540B \citep{minerva}. Furthermore, since the DeepSeekMath Corpus spans multiple languages, we observe significant gains on Chinese mathematical benchmark suites \citep{wei2023cmath,agieval}. We contend that our methodology in compiling and processing mathematical data establishes a valuable precedent for subsequent investigations, with ample prospects for future refinement.

DeepSeekMath-Base is built upon DeepSeek-Coder-Base-v1.5 7B \citep{deepseek-coder}, based on our observation that models pretrained on code perform better than those originating from general-purpose LLMs. In addition, our findings reveal that specialized math training elevates the model's performance not only in mathematical tasks but also on benchmarks like MMLU \citep{mmlu} and BBH \citep{bbh}, suggesting improvements in overall reasoning skills as well.

Following the pre-training phase, we fine-tune DeepSeekMath-Base using mathematical instruction datasets that incorporate chain-of-thought \citep{cot}, program-of-thought \citep{pot,pal}, and tool-integrated reasoning \citep{tora} approaches. This yields DeepSeekMath-Instruct 7B, which outperforms other 7B models and demonstrates performance on par with 70B open-source models trained with instruction tuning.

Additionally, we present Group Relative Policy Optimization (GRPO), a novel reinforcement learning (RL) approach derived from Proximal Policy Optimization (PPO) \citep{schulman2017proximal}. Unlike PPO, GRPO eliminates the need for a critic model by leveraging group scores to estimate the baseline, which considerably lowers computational demands during training. Employing only a portion of English instruction tuning data, GRPO demonstrates notable performance gains over the robust DeepSeekMath-Instruct, achieving improvements across both in-domain datasets (GSM8K: 82.9\% $\rightarrow$ 88.2\%, MATH: 46.8\% $\rightarrow$ 51.7\%) as well as out-of-domain mathematical benchmarks (for instance, CMATH: 84.6\% $\rightarrow$ 88.8\%) throughout the RL stage. We also articulate a coherent framework that unifies a range of methods, including Rejection Sampling Fine-Tuning (RFT) \citep{yuan2023scaling}, Direct Preference Optimization (DPO) \citep{dpo}, PPO, and GRPO. Within this unified framework, we observe that all these approaches can be interpreted as forms of direct or simplified RL methodologies. Our extensive experimental analysis examines key aspects of the paradigm—such as online versus offline training, supervision on outcome versus process, and comparisons between single-step and iterative RL setups—providing deeper insights into fundamental design choices. Finally, we discuss the mechanisms through which RL enhances the performance of instruction-tuned models and highlight promising avenues for future work towards more efficient RL, as guided by this unified perspective.

\subsection{Contributions}
We present scalable mathematical pre-training methods and further investigate reinforcement learning through comprehensive analysis.

\textbf{Math Pre-Training at Scale}

\begin{itemize}[topsep=0pt]
    \item This work demonstrates that the openly available Common Crawl dataset holds significant resources suitable for mathematical tasks. Through a carefully crafted data selection process, we assemble the DeepSeekMath Corpus, consisting of 120B tokens derived from web content filtered specifically for mathematical relevance. This dataset surpasses the size of math web pages employed by Minerva \citep{minerva} by nearly a factor of 7 and is approximately 9 times larger than the recently introduced OpenWebMath \citep{openwebmath}.

\item The DeepSeekMath-Base 7B pre-trained model demonstrates performance on par with Minerva 540B \citep{minerva}, suggesting that sheer parameter count is not the sole determinant of mathematical reasoning proficiency. High-quality pre-training enables smaller models to attain impressive results.

\item We present results from experiments focused on mathematical training. Conducting code training before engaging in math tasks enhances models’ performance on mathematical problem solving, regardless of whether auxiliary tools are involved. This partially addresses the enduring question: \textit{does code training enhance reasoning capabilities?} Our evidence suggests that, in the context of mathematical reasoning, it does.

\item Despite the widespread use of arXiv papers for training—particularly within mathematics-focused research—this approach yields no significant gains on any of the mathematical benchmarks evaluated in the present study.
\end{itemize}

\textbf{Exploration and Analysis of Reinforcement Learning}

\begin{itemize}[topsep=0pt]
    \item We propose Group Relative Policy Optimization (GRPO), a reinforcement learning method that operates without a critic model, leveraging group-based score estimation for the baseline. This approach leads to substantial savings in training resources in comparison to Proximal Policy Optimization (PPO).
    \item Our results show that GRPO considerably boosts the capability of our instruction-tuned model, DeepSeekMath-Instruct, relying exclusively on instruction-tuning data. Additionally, we detect improvements in generalization to out-of-domain tasks during reinforcement learning.
    \item We introduce a unified analytical framework that encapsulates various approaches such as RFT, DPO, PPO, and GRPO. A series of thorough experiments—covering topics like online versus offline training, outcome versus process supervision, and single-turn versus iterative reinforcement learning—are performed to rigorously analyze the foundational components of this framework.
    \item Drawing from our unified framework, we investigate underlying mechanisms that drive reinforcement learning's effectiveness, and identify promising avenues for enhancing reinforcement learning approaches for LLMs.
\end{itemize}

\subsection{Summary of Evaluations and Metrics}

\begin{itemize}[topsep=0pt]
    \item \textbf{English and Chinese Mathematical Reasoning}: Our analysis spans a variety of English and Chinese benchmarks designed to test mathematical reasoning skills at levels ranging from elementary school to university. For English tasks, we utilize benchmarks such as GSM8K \citep{gsm8k}, MATH \citep{MATH}, SAT \citep{llemma}, OCW Courses \citep{minerva}, and MMLU-STEM \citep{mmlu}. In Chinese, the evaluation encompasses MGSM-zh \citep{mgsm}, CMATH \citep{wei2023cmath}, Gaokao-MathCloze \citep{agieval}, and Gaokao-MathQA \citep{agieval}. Our evaluations investigate the models' proficiency in producing detailed, independent text-based solutions as well as their capabilities in resolving problems through the application of Python programming.

On English evaluation benchmarks, DeepSeekMath-Base demonstrates performance on par with the proprietary Minerva 540B model \citep{minerva}, and it outperforms all available open-source base models, such as Mistral 7B \citep{mistral} and Llemma-34B \citep{llemma}, regardless of whether these models have received mathematical pre-training. In many cases, this outperformance is by a wide margin. Furthermore, DeepSeekMath-Base achieves even stronger results on Chinese benchmarks, which can be attributed to its inclusion of high-quality mathematical data from multiple languages, rather than restricting pre-training solely to English as done in prior studies \citep{minerva,llemma}. Building on this foundation, applying mathematical instruction tuning and reinforcement learning yields DeepSeekMath-Instruct and DeepSeekMath-RL, models which achieve robust results—exceeding 50\% accuracy—on the competition-level MATH dataset, marking a first for open-source models.

\item \textbf{Formal Mathematics}: To assess DeepSeekMath-Base's capabilities, we utilize the informal-to-formal theorem proving benchmark as described in \citep{dsp_proof}, applied to the miniF2F dataset \citep{minif2f}, with Isabelle \citep{isabelle} serving as the chosen proof assistant. Results indicate that DeepSeekMath-Base achieves robust few-shot autoformalization outcomes.
\item \textbf{Natural Language Understanding, Reasoning, and Code}: In order to provide a thorough evaluation of the models' competencies in understanding, reasoning, and code generation, we test DeepSeekMath-Base on the Massive Multitask Language Understanding (MMLU) suite \citep{mmlu}, featuring 57 multiple-choice problems spanning a wide range of disciplines; on BIG-Bench Hard (BBH) \citep{bbh}, which includes 23 complex tasks emphasizing multi-step reasoning; and on code-related assessments such as HumanEval \citep{codex} and MBPP \citep{mbpp}, which are industry standards for benchmarking code models. Pre-training on mathematical data enhances both the model’s linguistic reasoning and its overall comprehension proficiency.

\section{Math Pre-Training}

\subsection{Data Collection and Decontamination} This section details the methodology used to build the DeepSeekMath Corpus utilizing Common Crawl data. As illustrated in Figure \ref{fig:data_collect}, we describe a cyclic pipeline designed for the organized extraction of a substantial mathematical dataset from Common Crawl, beginning with an initial seed corpus comprising a limited yet high-quality set of math-specific data. It is important to emphasize that this framework is adaptable and may be extended to other domains, including code.

Initially, we select OpenWebMath \citep{openwebmath}—a repository of high-quality mathematical content sourced from the web—as our seed dataset. Leveraging this collection, we develop a fastText model \citep{joulin2016fasttext} aimed at retrieving additional mathematical web pages with characteristics similar to those found in OpenWebMath. To train the model, we randomly sample 500,000 entries from the seed dataset as positive examples and another 500,000 web pages from Common Crawl to serve as negatives. The open-source fastText toolkit\footnote{\url{https://fasttext.cc}} is utilized for model training, with the following hyperparameters: a vector size of 256, learning rate set to 0.1, word n-gram lengths capped at 3, a minimum word frequency of 3, and 3 training epochs. To pare down the large Common Crawl dataset, we apply deduplication both by exact URL matches and via near-duplicate detection, ultimately producing a filtered corpus of 40B distinct HTML web pages. Employing the trained fastText model, we then extract web pages relevant to mathematics from this deduplicated corpus. To ensure quality, candidate pages are scored by the model, and only those with the highest predicted scores are retained. The final subset size is chosen based on performance in preliminary pre-training runs evaluating the top 40B, 80B, 120B, and 160B tokens. For the initial round, we opt to maintain the top 40B tokens.

Following the initial round of data acquisition, a substantial portion of mathematical web pages remains undiscovered. This is primarily attributed to the fastText model being trained on a set of positive samples lacking adequate diversity. To address this limitation, we identify additional mathematical web sources, thereby expanding the seed corpus in an effort to enhance the effectiveness of the fastText model. To do so, we partition the entirety of Common Crawl by domains, where each domain encompasses web pages with a common base URL. For each domain, we assess the proportion of pages obtained during the first iteration. Domains in which more than 10\% of their web pages are included in this initial collection are designated as math-related (for instance, \url{mathoverflow.net}).

Next, we proceed to manually label URLs within these domains that are indicative of mathematical content (such as \url{mathoverflow.net/questions}). Those pages linked to these URLs but absent from the original collection are then incorporated into the seed corpus. This strategy facilitates the accrual of more diverse positive samples, thus yielding an enhanced fastText model with improved recall for mathematical data in the following round of data collection. By the completion of four such iterations, a total of 35.5M mathematical web pages—amounting to 120B tokens—have been gathered. During the fourth iteration, it becomes apparent that approximately 98\% of the data was already present after the third iteration, prompting the cessation of further data collection.

To prevent benchmark contamination, we adopt the filtering procedure described in \cite{deepseek-coder}, excluding any web content that includes questions or answers derived from English mathematical benchmarks like GSM8K \citep{gsm8k}, MATH \citep{MATH}, or from Chinese benchmarks such as CMATH \citep{wei2023cmath} and AGIEval \citep{agieval}. Specifically, we remove from our math training dataset any text segment containing an exact match to any 10-gram phrase present in the evaluation benchmarks. In the case of benchmark passages shorter than 10 grams but consisting of at least 3 grams, we apply exact-match filtering to eliminate web pages containing those sequences.

\subsection{Validating the Quality of the DeepSeekMath~Corpus}

We run pre-training experiments to investigate how the DeepSeekMath~Corpus is compared with the recently released math-training corpora:

\begin{itemize}[topsep=0pt]
    \item \textbf{MathPile} \citep{mathpile}: This is a composite dataset comprising 8.9B tokens, created by merging resources such as textbooks, Wikipedia, ProofWiki, CommonCrawl, StackExchange, and arXiv, with arXiv providing the predominant share (over 85\%) of its contents;
    \item \textbf{OpenWebMath} \citep{openwebmath}: Derived from CommonCrawl, this resource selects specifically for mathematical material, resulting in a total of 13.6B tokens;
    \item \textbf{Proof-Pile-2} \citep{llemma}: This dataset integrates OpenWebMath, AlgebraicStack (which contributes 10.3B tokens from mathematical code), and arXiv documents (contributing 28.0B tokens). In experiments involving Proof-Pile-2, we adhere to the arXiv:Web:Code proportion of 2:4:1, as specified by \cite{llemma}.
\end{itemize}

\subsubsection{Training Setting}
\label{sec:quality-policy}
We conduct math-specific training on a foundational language model comprising 1.3B parameters, architecturally aligned with the DeepSeek LLMs \citep{deepseek-llm}, herein referenced as DeepSeek-LLM 1.3B. For each distinct mathematical dataset, a separate model is trained over 150B tokens. All training procedures utilize the streamlined and resource-efficient HAI-LLM framework \citep{haillm}. Consistent with the established protocols for DeepSeek LLMs, the AdamW optimizer \citep{adamW} is employed with parameter settings of $\beta_1=0.9$, $\beta_2=0.95$, and $\mathrm{weight\_decay}=0.1$. A multi-phase learning rate schedule is implemented: the learning rate ramps up to its maximum over the first 2,000 warmup steps, declines to 31.6\% of its peak by the 80\% completion mark, and further reduces to 10.0\% of the peak at the 90\% completion point. The learning rate’s upper bound is fixed at 5.3e-4, with training utilizing a batch size of 4M tokens and a 4K token context window.

\subsubsection{Evaluation Results}

\textbf{The DeepSeekMath~Corpus is of high quality, covers multilingual mathematical content, and is the largest in size.}

\begin{itemize}[topsep=0pt]
    \item \textbf{High-quality}:
    We assess the models’ downstream capabilities on eight mathematical evaluation sets utilizing few-shot chain-of-thought prompting \cite{cot}. According to Table \ref{tab:corpora_comparison}, models trained with the DeepSeekMath~Corpus exhibit a marked advantage in performance. Figure \ref{fig:corpora_comparisons} further illustrates that the model trained with the DeepSeekMath Corpus surpasses Proof-Pile-2 at the 50B token mark (corresponding to a single complete pass through Proof-Pile-2), reflecting the superior average quality of the DeepSeekMath Corpus.
    \item \textbf{Multilingual}:
    The DeepSeekMath~Corpus includes content in a range of languages, with English and Chinese being the most prominent. Table \ref{tab:corpora_comparison} shows that utilizing the DeepSeekMath~Corpus for training yields improved mathematical reasoning proficiency in both English and Chinese. By contrast, currently available mathematical corpora are predominantly English-focused and contribute minimally—or can even adversely affect—performance in Chinese mathematical reasoning tasks.
    \item \textbf{Large-scale}:
    The size of the DeepSeekMath~Corpus exceeds that of previous mathematical corpora by a significant margin. Figure \ref{fig:corpora_comparisons} shows that DeepSeek-LLM 1.3B, trained on the DeepSeekMath~Corpus, benefits from a much sharper learning progression and achieves longer-lasting gains. In contrast, the smaller baseline corpora must be cycled through multiple times in training, causing model improvement to stagnate at an earlier stage.
\end{itemize}

\subsection{Training and Evaluating DeepSeekMath-Base 7B}

This section presents DeepSeekMath-Base 7B, a foundational model exhibiting notable proficiency in reasoning tasks, particularly within mathematical domains. The model is based on DeepSeek-Coder-Base-v1.5 7B \citep{deepseek-coder} and undergoes training on a dataset comprising 500B tokens. The dataset composition is distributed as follows: 56\% from the DeepSeekMath Corpus, 4\% from AlgebraicStack, 10\% drawn from arXiv, 20\% consisting of code from Github, and the remaining 10\% consisting of natural language text sourced from Common Crawl in English and Chinese. Training largely follows the configuration described in Section \ref{sec:quality-policy}, with the exception of assigning the peak learning rate to 4.2e-4 and utilizing a batch size totaling 10M tokens.

We systematically evaluate the mathematical proficiency of DeepSeekMath-Base 7B by examining its competence in independently generating complete mathematical solutions, utilizing external tools to solve mathematical tasks, and performing rigorous formal theorem verification. Additionally, we extend our analysis to capture a broader overview of the base model’s general capabilities, highlighting its effectiveness in natural language comprehension, logical reasoning, and programming tasks.

\paragraph{Mathematical Problem Solving with Step-by-Step Reasoning}

We assess the capabilities of DeepSeekMath-Base in mathematical problem solving via few-shot chain-of-thought prompting \citep{cot}, utilizing eight distinct benchmarks in both English and Chinese. These benchmarks span quantitative reasoning tasks—such as those found in GSM8K \citep{gsm8k}, MATH \citep{MATH}, and CMATH \citep{wei2023cmath}—as well as multiple-choice formats, exemplified by MMLU-STEM \citep{mmlu} and Gaokao-MathQA \citep{agieval}. Collectively, they address a broad spectrum of mathematical topics ranging from elementary arithmetic to advanced college-level challenges.

According to Table \ref{tab:base_math_cot}, DeepSeekMath-Base 7B demonstrates superior results across all eight benchmarks when compared to other open-source base models, such as the prominent general-purpose model Mistral 7B \citep{mistral} and the more recent Llemma 34B \citep{llemma}, the latter having received additional mathematical training on Proof-Pile-2 \citep{llemma}. Especially noteworthy is its performance on the MATH dataset at competition level, where DeepSeekMath-Base exceeds all open-source base models by more than 10 percentage points, and even outperforms Minerva 540B \citep{minerva}—a proprietary base model with 77 times the parameter count, built upon PaLM \citep{palm} and supplemented with further mathematical corpus training.

\paragraph{Mathematical Problem Solving with Tool Use} We assess mathematical reasoning that incorporates external tools on the GSM8K and MATH benchmarks, employing few-shot program-of-thought prompting methods \citep{pot,pal}. For each problem, the models are instructed to generate a Python script, leveraging standard libraries like \textit{math} and \textit{sympy} for advanced calculations as needed. The answer is determined by executing the generated program. According to Table \ref{tab:base_math_pot_proof}, DeepSeekMath-Base 7B achieves superior performance compared to the previous leading model, Llemma 34B.

\paragraph{Formal Mathematics}

Automating formal proofs offers significant advantages in terms of ensuring correctness and robustness of mathematical arguments, as well as improving workflow efficiency—a topic that has gained considerable interest in recent years. In this work, we assess the capabilities of DeepSeekMath-Base 7B for the informal-to-formal proving task presented in \citep{dsp_proof}, wherein the goal is to produce a formalized proof given an informal proposition, its formalized version, and an informal demonstration. Our evaluation employs the miniF2F dataset \citep{minif2f}, a standard benchmark for high-level mathematical formalization, requiring the generation of formal Isabelle proofs for each problem using a few-shot prompting approach. Adhering to the methodology in \cite{dsp_proof}, we have the model generate high-level proof sketches, after which we utilize the Sledgehammer automated theorem prover \citep{sledgehammer} to elaborate and complete the formal details. Results in Table \ref{tab:base_math_pot_proof} highlight that DeepSeekMath-Base 7B achieves competitive performance in automated proof formalization tasks.

\paragraph{Natural Language Understanding, Reasoning, and Code}

We assess natural language understanding using MMLU \citep{mmlu}, evaluate reasoning abilities via BBH \citep{bbh}, and measure coding proficiency with HumanEval \citep{codex} and MBPP \citep{mbpp}. According to Table \ref{tab:understanding_reasoning_code}, DeepSeekMath-Base 7B achieves notable improvements over its predecessor, DeepSeek-Coder-Base-v1.5 \citep{deepseek-coder}, in both MMLU and BBH tasks. This demonstrates that specialized math training benefits language comprehension and reasoning competencies. Furthermore, by incorporating code tokens during further training, DeepSeekMath-Base 7B is able to sustain the strong coding performance observed in DeepSeek-Coder-Base-v1.5 on both coding evaluations. In summary, DeepSeekMath-Base 7B yields a substantial advantage compared to the generalist Mistral 7B \citep{mistral} model on all three reasoning and coding assessments.

\section{Supervised Fine-Tuning}

\subsection{SFT Data Curation}

We develop a comprehensive instruction-tuning dataset for mathematics that includes both English and Chinese questions across a range of mathematical disciplines and difficulty levels. Each question is matched with an accompanying solution presented in multiple reasoning formats: chain-of-thought (CoT) \citep{cot}, program-of-thought (PoT) \citep{pot,pal}, and tool-integrated reasoning \citep{tora}. In total, the dataset consists of 776K training samples.

\begin{itemize}[topsep=0pt]
    \item \textbf{English mathematical datasets}: We provide tool-integrated solution annotations for the GSM8K and MATH datasets, and incorporate selected items from MathInstruct \citep{MathInstruct} along with the Lila-OOD training split \citep{lila}, both of which use either CoT or PoT approaches to problem solving. This English dataset set encompasses a wide array of mathematical disciplines, including calculus, number theory, algebra, probability, and geometry.
    \item \textbf{Chinese mathematical datasets}: We assemble a set of K-12 mathematics problems in Chinese, covering 76 distinct subtopics such as linear equations. Each problem is annotated with both CoT-based and tool-supported solution strategies.
\end{itemize}

\subsection{Training and Evaluating DeepSeekMath-Instruct 7B}

This section presents DeepSeekMath-Instruct 7B, a model refined with mathematical instruction tuning using DeepSeekMath-Base as its foundation. For training, examples are randomly joined together until the sequence attains a maximum context length of 4K tokens. The model is trained across 500 steps, employing a batch size of 256 and maintaining a fixed learning rate of 5e-5 throughout.

We evaluate models' mathematical performance both without and with tool use, on 4 quantitative reasoning benchmarks in English and Chinese. We benchmark our model against the leading models of the time:

\begin{itemize}[topsep=0pt]
    \item \textbf{Closed-source models} are represented by: (1) the GPT series, notably GPT-4 \citep{gpt4} and GPT-4 Code Interpreter~\footnote{\url{https://openai.com/blog/chatgpt-plugins\#\#code-interpreter}}, recognized as the most advanced in this group, (2) Gemini Ultra and Pro \citep{gemini}, (3) Inflection-2 \citep{inflection-2}, (4) Grok-1~\footnote{\url{https://x.ai/model-card}}, and a selection of newly released models from Chinese technology firms including (5) Baichuan-3~\footnote{\url{https://www.baichuan-ai.com}}, and (6) the newest GLM-4~\footnote{\url{https://open.bigmodel.cn/dev/api\#glm-4}}

    from the GLM family \citep{glm}.

These models are designed for general use cases, and most have been subjected to various alignment methods.

\item \textbf{Open-source models} encompass: broad-coverage systems such as (1) DeepSeek-LLM-Chat 67B \citep{deepseek-llm}, (2) Qwen 72B \citep{qwen}, (3) SeaLLM-v2 7B \citep{seallm}, and (4) ChatGLM3 6B \citep{chatglm3}. Additionally, there are specialized models with improved mathematical capabilities, including: (5) InternLM2-Math 20B~\footnote{\url{https://github.com/InternLM/InternLM-Math}}, developed from InternLM2 and further trained on mathematics followed by instruction tuning; (6) Math-Shepherd-Mistral 7B, which incorporates PPO training \citep{schulman2017proximal} on Mistral 7B \citep{mistral} guided by a process-supervised reward model; (7) the WizardMath family \citep{wizardmath}, which leverages evolve-instruct (a variant of instruction tuning utilizing AI-generated instructions) and PPO training to strengthen mathematical reasoning in Mistral 7B and Llama-2 70B \citep{llama2}, with datasets drawn mainly from GSM8K and MATH; (8) MetaMath 70B \citep{metamath}, which adapts Llama-2 70B through fine-tuning on enriched versions of GSM8K and MATH; (9) ToRA 34B \cite{tora}, built by fine-tuning CodeLlama 34B for tool-integrated mathematical problem solving; and (10) MAmmoTH 70B \citep{MathInstruct}, based on Llama-2 70B with additional instruction tuning utilizing MathInstruct.

According to the results in Table \ref{tab:sft_rl_math}, when evaluated in a scenario that prohibits the use of external tools, DeepSeekMath-Instruct 7B achieves notable proficiency in step-by-step problem-solving. Specifically, for the challenging competition-tier MATH benchmark, our model outperforms every open-source alternative and most closed-source systems (including Inflection-2 and Gemini Pro) by a margin of at least 9 percentage points. This remains the case even in comparison to considerably larger models (such as Qwen 72B) or those benefiting from targeted mathematical reinforcement learning strategies (e.g., WizardMath-v1.1 7B). Although DeepSeekMath-Instruct attains results comparable to proprietary Chinese models like GLM-4 and Baichuan-3 on the MATH dataset, it is still outperformed by GPT-4 and Gemini Ultra.

When assessed in an environment that permits the combination of natural language reasoning with programmatic tool usage for solving problems, DeepSeekMath-Instruct 7B achieves nearly 60\% accuracy on the MATH benchmark, outperforming every other open-source model available to date. For additional benchmarks, our model delivers results comparable to those of DeepSeek-LLM-Chat 67B—the former leading model in this domain—despite its own parameter count being only one-tenth as large.
\section{Reinforcement Learning}

\subsection{Group Relative Policy Optimization} Recent studies have demonstrated that reinforcement learning (RL) serves as a valuable tool for enhancing the mathematical reasoning capabilities of LLMs following Supervised Fine-Tuning (SFT) \citep{wang2023math,wizardmath}. Here, we present Group Relative Policy Optimization (GRPO), an RL method designed for both efficiency and effectiveness.

\subsubsection{From PPO to GRPO} Proximal Policy Optimization (PPO) \citep{schulman2017proximal} is an actor-critic RL algorithm that is widely used in the RL fine-tuning stage of LLMs \citep{ouyang2022training}. In particular, it optimizes LLMs by maximizing the following surrogate objective:

\begin{equation}
\footnotesize
    \mathcal{J}_{PPO}(\theta) = \mathbb{E}{[q \sim P(Q), o \sim \pi_{\theta_{old}}(O|q)]} \frac{1}{|o|} \sum_{t=1}^{|o|} \min \left[ \frac{\pi_\theta(o_{t} | q, o_{<t})}{\pi_{\theta_{old}}(o_{t} | q, o_{<t})} A_{t}, \text{clip} \left( \frac{\pi_\theta(o_{t} | q, o_{<t})}{\pi_{\theta_{old}}(o_{t} | q, o_{<t})}, 1 - \varepsilon, 1 + \varepsilon \right)  A_{t} \right] ,
\end{equation}

The above loss function for Proximal Policy Optimization is defined as the expectation, with respect to $q$ drawn from the distribution $P(Q)$ and $o$ sampled according to the previous policy $\pi_{\theta_{old}}(O|q)$, of the averaged minimum over time steps $t$ in each output sequence. At each step, it computes the minimum between the advantage-weighted likelihood ratio for the current and previous policies, and a clipped version of the same ratio (bounded by $1-\varepsilon$ and $1+\varepsilon$), both scaled by the corresponding advantage $A_t$.

In this context, $\pi_{\theta}$ and $\pi_{\theta_{old}}$ denote the present and previous iterations of the policy models, while $q$ and $o$ refer to the questions drawn from the question dataset and outputs generated by the prior policy $\pi_{\theta_{old}}$, respectively. The parameter $\varepsilon$ serves as a clipping hyper-parameter within PPO to ensure greater training stability. The term $A_t$ represents the advantage, calculated through the Generalized Advantage Estimation (GAE) technique \citep{gae}, utilizing the subsequent rewards $\{r_{\ge t}\}$ together with an estimated value function $V_{\psi}$. Consequently, PPO requires simultaneous optimization of both the value function and the policy model. To address potential issues of reward model over-optimization, a KL divergence penalty with respect to a reference model is commonly integrated into the per-token reward during training, as detailed by \citet{ouyang2022training}, that is,

\begin{equation}
    r_{t} = r_\varphi(q, o_{\le t}) - \beta \log\left(\frac{\pi_{\theta}(o_{t}|q, o_{<t})}{\pi_{ref}(o_{t}|q, o_{<t})}\right),
\label{eq:PPO-reward}
\end{equation}

In this context, $r_\varphi$ denotes the reward model, while $\pi_{ref}$ refers to the reference model—typically represented by the original SFT model. The parameter $\beta$ serves as the weighting factor for the KL divergence penalty.

As the value function employed in PPO is typically another model of comparable size as the policy model, it brings a substantial memory and computational burden. Additionally, during RL training, the value function is treated as a baseline in the calculation of the advantage for variance reduction. While in the LLM context, usually only the last token is assigned a reward score by the reward model, which may complicate the training of a value function that is accurate at each token. To address this, as shown in Figure \ref{fig:grpo}, we propose Group Relative Policy Optimization (GRPO), which obviates the need for additional value function approximation as in PPO, and instead uses the average reward of multiple sampled outputs, produced in response to the same question, as the baseline. More specifically, for each question $q$, GRPO samples a group of outputs $\{o_1, o_2, \cdots, o_G\}$  from the old policy  $\pi_{\theta_{old}}$  and then optimizes the policy model by maximizing the following objective:

\begin{equation}
\footnotesize
\begin{split}
    \mathcal{J}_{GRPO}(\theta) &= \mathbb{E}_{q \sim P(Q),\; \{o_i\}_{i=1}^G \sim \pi_{\theta_{old}}(O|q)}  \\
    & \frac{1}{G}\sum_{i=1}^G \frac{1}{|o_i|} \sum_{t=1}^{|o_i|} \Bigg[ \min \Bigg( \frac{\pi_\theta(o_{i,t} | q, o_{i,<t})}{\pi_{\theta_{old}}(o_{i,t} | q, o_{i,<t})} \hat{A}_{i,t}, \\
    & \hspace{5em} \text{clip} \left( \frac{\pi_\theta(o_{i,t} | q, o_{i,<t})}{\pi_{\theta_{old}}(o_{i,t} | q, o_{i,<t})}, 1 - \varepsilon, 1 + \varepsilon \right)  \hat{A}_{i,t} \Bigg) - \beta \mathbb{D}_{KL}\left[\pi_{\theta} || \pi_{ref}\right] \Bigg],
\end{split}
\label{eq:GRPO-obj}
\end{equation}

where $\varepsilon$ and $\beta$ are hyper-parameters, and $\hat{A}_{i,t}$  is the advantage calculated based on relative rewards of the outputs inside each group only, which will be detailed in the following subsections. The group relative way that GRPO leverages to calculate the advantages, aligns well with the comparative nature of rewards models,  as reward models are typically trained on datasets of comparisons between outputs on the same question. Also note that, instead of adding KL penalty in the reward, GRPO regularizes by directly adding the KL divergence between the trained policy and the reference policy to the loss, avoiding complicating the calculation of  $\hat{A}_{i,t}$. And different from the KL penalty term used in (\ref{eq:PPO-reward}), we estimate the KL divergence with the following unbiased estimator \citep{kl_approx}:

\begin{equation}
\small
    \mathbb{D}_{KL}\left[\pi_{\theta} || \pi_{ref}\right] = \frac{\pi_{ref}(o_{i,t}|q,o_{i,<t})}{\pi_{\theta}(o_{i,t}|q,o_{i,<t})} - \log\left(\frac{\pi_{ref}(o_{i,t}|q,o_{i,<t})}{\pi_{\theta}(o_{i,t}|q,o_{i,<t})}\right) - 1,
\end{equation}

which is guaranteed to be positive.

\begin{algorithm}[t]
  \small
  \caption{Iterative Group Relative Policy Optimization}
  \textbf{Input} initial policy network $\pi_{\theta_{\text{init}}}$; collection of reward models $r_\varphi$; set of task prompts $\mathcal{D}$; 
  hyperparameters $\varepsilon$, $\beta$, $\mu$
  \begin{algorithmic}[1]
    \State Set policy network $\pi_\theta \leftarrow \pi_{\theta_{\text{init}}}$
    \For{iteration = 1, \dots, I}
       \State Assign $\pi_{ref} \leftarrow \pi_{\theta}$ as the reference model
      \For{step = 1, \dots, M}
      \State Randomly select a minibatch $\mathcal{D}_b$ from $\mathcal{D}$
      \State Store current parameters: $\pi_{\theta_{old}} \leftarrow \pi_{\theta}$ 
      \State For each question $q$ in $\mathcal{D}_b$, generate $G$ candidate responses $\{o_i\}_{i=1}^G \sim \pi_{\theta_{old}} (\cdot \mid q) $
      \State Obtain corresponding rewards $\{r_i\}_{i=1}^{G}$ for every $o_i$ by applying $r_{\varphi}$
      \State Estimate $\hat{A}_{i,t}$ for token $t$ of each $o_i$ using group-relative advantage computation
      \For{GRPO iteration = 1, \dots, $\mu$}
        \State Update $\pi_{\theta}$ by maximizing the group relative policy optimization objective (Equation \ref{eq:GC-GRPO})
      \EndFor
    \EndFor
    \State Refine reward models $r_\varphi$ via continual learning, leveraging experience replay.
    \EndFor 
  \end{algorithmic}
  \textbf{Output} $\pi_\theta$
  \label{alg:iter-grpo}
\end{algorithm}

\subsubsection{Outcome Supervision RL with GRPO} Specifically, given a question $q$, a set of outputs $\{o_1, o_2, \cdots, o_G\}$ is drawn from the previous policy model $\pi_{\theta_{old}}$. Each output is evaluated using a reward model, which generates a collection of $G$ reward values $\mathbf{r} = \{r_1, r_2, \cdots, r_G\}$. These reward values are then standardized by removing the group mean and scaling by the group standard deviation. Under outcome supervision, the normalized reward for each response $o_i$ is provided at the completion of the output. The advantage estimate $\hat{A}_{i, t}$ for every token in $o_i$ is assigned to this normalized reward, that is, $\hat{A}_{i, t} = \widetilde{r}_i = \frac{r_i- {\rm mean}(\mathbf{r})}{{\rm std}(\mathbf{r})}$. The policy is subsequently updated by maximizing the objective shown in equation (\ref{eq:GRPO-obj}).

\subsubsection{Process Supervision RL with GRPO} While outcome supervision delivers a single reward only at the completion of each output, this feedback can be insufficient and inefficient for guiding the policy on intricate mathematical problems. Building upon \cite{wang2023math}, we also investigate process supervision, in which feedback is given after each individual reasoning step. Specifically, for a given question $q$ and a collection of $G$ sampled outputs $\{o_1, o_2, \cdots, o_G\}$, a process reward model evaluates every step within the outputs and assigns corresponding rewards: $\mathbf{R} = \{ \{r_1^{index(1)},\cdots,r_1^{index(K_1)}\}, \cdots,  \{r_G^{index(1)},\cdots,r_G^{index(K_G)}\} \}$. Here, $index(j)$ denotes the end-of-step token index for the $j$-th reasoning step, and $K_i$ specifies the total number of reasoning steps for the $i$-th output. These stepwise rewards are standardized by subtracting their mean and dividing by the standard deviation: $\widetilde{r}_i^{index(j)} = \frac{r_i^{index(j)} - {\rm mean(\mathbf{R})}}{{\rm std(\mathbf{R})}}$.

Process supervision proceeds by computing the token-level advantage as the cumulative sum of the normalized rewards assigned to all subsequent steps, that is, $\hat{A}_{i, t} = \sum_{index(j) \ge t} \widetilde{r}_i^{index(j)}$. The policy is then updated by maximizing the objective presented in equation (\ref{eq:GRPO-obj}).

\subsubsection{Iterative RL with GRPO} During the course of reinforcement learning, the existing reward model may become inadequate for effectively guiding the evolving policy model. To address this, we investigate iterative RL with GRPO. As illustrated in Algorithm \ref{alg:iter-grpo}, this approach involves generating fresh datasets for the reward model by sampling outputs from the current policy model, and subsequently updating the reward model through a replay strategy that retains 10\% of previous data. Afterward, the reference model is updated to match the current policy model, and the policy model itself is further optimized using the improved reward model.

\subsection{Training and Evaluating DeepSeekMath-RL}

Our reinforcement learning experiments utilize the DeepSeekMath-Instruct 7B model. The RL training data consists exclusively of chain-of-thought-style problems from GSM8K and MATH present in the SFT corpus, totaling approximately 144,000 questions. We deliberately exclude all other SFT tasks in order to more precisely observe the effect of RL on benchmarks that are not represented during the RL process. The reward model's training dataset is built in accordance with the methodology of \citep{wang2023math}. For reward model training, we employ DeepSeekMath-Base 7B as the backbone, using a learning rate of $2 \times 10^{-5}$. For the policy update under GRPO, the learning rate is set to $1 \times 10^{-6}$, and we use a KL divergence coefficient of 0.04. Each question receives $64$ sampled completions. Both the maximum sequence length and training batch size are fixed at 1024. After each exploration phase, the policy undergoes a single update step. In evaluation, DeepSeekMath-RL 7B is benchmarked following the protocols established by DeepSeekMath-Instruct 7B. Within this experimental setup, chain-of-thought GSM8K and MATH are treated as in-domain tasks, while all other test sets are considered out-of-domain.

Table \ref{tab:sft_rl_math} presents a comparative analysis of open-source and closed-source models employing both chain-of-thought and tool-based reasoning approaches on English and Chinese evaluation datasets. Our findings include: 1) DeepSeekMath-RL 7B achieves accuracy rates of 88.2\% on GSM8K and 51.7\% on MATH when using chain-of-thought reasoning. These results outperform all open-source models within the 7B to 70B scale, as well as most closed-source counterparts. 2) Notably, DeepSeekMath-RL 7B is exclusively fine-tuned on chain-of-thought instruction data from GSM8K and MATH, originating from DeepSeekMath-Instruct 7B. In spite of the limited diversity of its training data, it consistently exceeds the performance of DeepSeekMath-Instruct 7B on all evaluation criteria, highlighting the advantages conferred by reinforcement learning.

\section{Discussion}
Here, we present an overview of the results obtained from both pre-training and reinforcement learning experiments.

\subsection{Lessons Learnt in Pre-Training}

We begin by detailing our observations from the pre-training phase. Except where stated, our methods correspond to the training configuration specified in Section \ref{sec:quality-policy}. Importantly, references to the DeepSeekMath Corpus in this discussion pertain to the 89-billion-token dataset obtained from the second cycle of data acquisition.

\subsubsection{Code Training Benefits Mathematical Reasoning}

An often-cited but unproven claim posits that training on code enhances reasoning abilities. Here, we seek to provide partial evidence addressing this claim, focusing on mathematics: code training augments models’ mathematical reasoning skills, regardless of whether tools are employed.

To study how code training affects mathematical reasoning, we experimented with the following two-stage training and one-stage training settings:

\textbf{Two-Stage Training}

\begin{itemize}[topsep=0pt]
    \item \textbf{Code Training for 400B Tokens $\rightarrow$ Math Training for 150B Tokens}: Our procedure first involves exposing DeepSeek-LLM 1.3B to 400B code tokens before transitioning to an additional 150B math token training phase;
    \item \textbf{General Training for 400B Tokens $\rightarrow$ Math Training for 150B Tokens}: For comparison, we perform a similar training regimen where, in the initial stage, general tokens drawn from a comprehensive general-purpose dataset curated by DeepSeek-AI replace code tokens. This setup allows us to assess whether pre-training on code confers any distinct advantage over general language pre-training in boosting mathematical reasoning abilities.
\end{itemize}

\textbf{One-Stage Training}

\begin{itemize}[topsep=0pt]
    \item \textbf{Math Training for 150B Tokens}: DeepSeek-LLM 1.3B undergoes training on 150B tokens specific to mathematics;
    \item \textbf{Training on a mixture of 400B Code Tokens and 150B Math Tokens}: Sequential training on math data after coding data diminishes coding abilities. We assess if concurrently training on a combined dataset comprising both code (400B tokens) and math (150B tokens) can enhance mathematical reasoning while also mitigating the issue of catastrophic forgetting observed in staged training.
\end{itemize}

\paragraph{Results}
The outcomes of downstream evaluations across various training configurations are presented in Table \ref{tab:code-math} and Table \ref{tab:code-math-general-eval}.

Code training serves to improve program-aided mathematical reasoning in both the two-stage and one-stage training paradigms. As indicated in Table \ref{tab:code-math}, in the context of two-stage training, engaging in code training alone markedly boosts performance on GSM8K and MATH tasks using Python, with additional gains achieved when math training is introduced in the subsequent stage. Notably, when applying one-stage training, incorporating a mix of code and math tokens helps prevent catastrophic forgetting observed in two-stage processes, and further enhances both general coding capabilities (Table \ref{tab:code-math-general-eval}) and program-aided mathematical reasoning (Table \ref{tab:code-math}).

Training with code similarly enhances mathematical reasoning skills even in the absence of tool assistance. In the two-stage training approach, introducing code training at the outset yields measurable gains in performance. This stage additionally facilitates more effective math training thereafter, culminating in superior overall results. In contrast, a single-stage training regime that merges code and mathematical tokens detracts from mathematical reasoning capability when tools are not involved. It is hypothesized that the relatively small parameter size of DeepSeek-LLM 1.3B restricts its ability to comprehensively integrate and process both code and mathematical content in a unified phase.

\subsubsection{ArXiv Papers Seem Ineffective in Improving Mathematical Reasoning}

ArXiv papers are commonly included as a component of math pre-training data \citep{minerva,gpt-f,llemma,mathpile}. However, detailed analysis regarding their impact on mathematical reasoning has not been extensively conducted. Perhaps counter-intuitively, according to our experiments, arXiv papers seem ineffective in improving mathematical reasoning. We experiment with models of different sizes, including DeepSeek-LLM 1.3B and DeepSeek-Coder-Base-v1.5 7B \citep{deepseek-coder}, using arXiv corpora that underwent varied processing pipelines:

\begin{itemize}[topsep=0pt]
    \item \textbf{MathPile} \citep{mathpile}: comprises 8.9 billion tokens, assembled using a set of cleaning and filtering heuristics, with more than 85\% of the dataset originating from scientific documents on arXiv;
    \item \textbf{ArXiv-RedPajama} \citep{redpajama}: contains all arXiv LaTeX source files, systematically stripped of preambles, comments, macros, and bibliographies, amounting to 28.0 billion tokens.
\end{itemize}

For our study, we independently pretrain DeepSeek-LLM 1.3B for 150B tokens and DeepSeek-Coder-Base-v1.5 7B for 40B tokens using the respective arXiv corpus. Results suggest that exposure to arXiv papers alone does not enhance models’ mathematical reasoning capabilities. Training solely on arXiv material fails to yield meaningful gains and in some cases leads to a decline in performance across a range of mathematical benchmarks of varying difficulty featured in this work. Among these assessments are quantitative reasoning tasks such as GSM8K and MATH (refer to Table \ref{tab:arxiv-cot}), multiple-choice tests including MMLU-STEM (Table \ref{tab:arxiv-cot}), as well as formal mathematics datasets like miniF2F (Table \ref{tab:arxiv_atp}).

However, this conclusion has its limitations and should be taken with a grain of salt. We have not yet studied:

\begin{itemize}[topsep=0pt]
    \item How arXiv tokens influence particular math-focused tasks omitted from this study, for example, transforming formal theorems or proofs into their informal counterparts (informalization);
    \item The interaction between arXiv tokens and additional data sources;
    \item If scaling up model size would reveal the advantages of utilizing arXiv papers.
\end{itemize}

Thus, further exploration is required, which we leave for future studies.

\subsection{Insights of Reinforcement Learning}
\subsubsection{Towards to a Unified Paradigm} In this part, we introduce a consolidated framework for examining various training strategies, including SFT, RFT, DPO, PPO, GRPO, and also carry out experiments to investigate the elements underlying this unified approach. Broadly, the gradient concerning the parameter $\theta$ for any training method may be expressed as:

\begin{equation}
    \nabla_{\theta}\mathcal{J}_{\textcolor{red}{\mathcal{A}}}(\theta) = \mathbb{E}[\underbrace{(q,o) \sim \textcolor{red}{\mathcal{D}}}_{Data \ Source}]\left( \frac{1}{|o|} \sum_{t=1}^{|o|}  \underbrace{GC_{{\mathcal{A}}}(q, o, t, \textcolor{red}{\pi_{{rf}}})}_{Gradient \ Coefficient}  \nabla_{\theta}\log \pi_{\theta}(o_t | q, o_{<t})\right).
\label{eq:objective}
\end{equation}

There exist three key components: 1) \textit{Data Source $\mathcal{D}$}, which determines the training data; 2) \textit{Reward Function $\pi_{{rf}}$}, which is the source of the training reward signal; 3) \textit{Algorithm $\mathcal{A}$}: which processes the training data and the reward signal to the gradient coefficient $GC$ that determines the magnitude of the penalty or reinforcement for the data. We analyze several representative methods based on such a unified paradigm:

\begin{itemize}
    \item  \textbf{Supervised Fine-tuning (SFT)}: SFT involves adjusting a pretrained model by training it on a dataset curated by humans for supervised learning.
    \item \textbf{Rejection Sampling Fine-tuning (RFT)}: RFT continues to train the SFT model, but exclusively on outputs generated and sampled from the SFT model itself in response to SFT questions, retaining only those whose answers are evaluated as correct.
    \item \textbf{Direct Preference Optimization (DPO)}: DPO further improves the SFT model by fine-tuning with a set of additional outputs created by the SFT model, where optimization is performed with respect to the DPO loss over paired preferences.
    \item \textbf{Online Rejection Sampling Fine-tuning (Online RFT)}: In contrast to RFT, Online RFT starts with the SFT model as the initial policy, then progressively enhances this policy by fine-tuning on newly generated samples drawn from the current version of the policy model.
    \item \textbf{PPO/GRPO}: PPO/GRPO begins with the SFT model for initialization and employs reinforcement learning to update the policy using outputs dynamically generated by the evolving policy model itself.
\end{itemize}

Table \ref{tab:data_policy} presents an overview of the key elements associated with these methods. For a comprehensive explanation of the derivation procedure, consult Appendix \ref{app:analysis-rl}.

\paragraph{Observation about Data Source} The data sources considered fall into two types: online sampling and offline sampling. Online sampling refers to data collected through the interactions of the current policy during training, whereas offline sampling refers to data obtained from samples generated by the original SFT model. Methods such as RFT and DPO utilize offline data, in contrast to Online RFT and GRPO, which rely on data acquired via online sampling.

As illustrated in Figure \ref{fig:rl-analysis}, our analysis reveals that Online RFT achieves markedly better performance than RFT across both benchmarks. At the onset of training, Online RFT and RFT display similar outcomes; however, as training advances, Online RFT establishes a clear lead, underscoring the benefits of online training strategies. This observation aligns with expectations: initially, the actor and SFT model remain quite similar, resulting in sampled data with only slight discrepancies. As training progresses, greater divergence emerges in the actor’s sampled data, thereby amplifying the benefits of collecting data in real-time.

\paragraph{Observation about Gradient Coefficient} In the update of model parameters, the algorithm utilizes the processed reward signal as the gradient coefficient. Within our experiments, the reward function is categorized into `Rule' and `Model'. The `Rule' approach assesses response quality strictly by evaluating answer correctness, whereas the `Model' approach relies on a reward model trained to assign scores to responses, with its training data originating from rule-based assessments. As reflected in Equations \ref{eq:GC-RFT} and \ref{eq:GC-GRPO}, GRPO and Online RFT diverge fundamentally: GRPO exclusively adapts the gradient coefficient dynamically according to the reward values output by the reward model. This adaptation makes it possible to strengthen or weaken responses proportionally based on the actual reward magnitude. In comparison, Online RFT does not incorporate this mechanism—incorrect responses are not penalized, and all correct responses are reinforced equally, irrespective of their underlying quality as reflected by the reward signal.

As illustrated in Figure \ref{fig:rl-analysis}, GRPO outperforms online RFT, emphasizing the advantage of modifying positive and negative gradient coefficients. Moreover, GRPO+PS achieves better results than GRPO+OS, which underscores the effectiveness of employing fine-grained, step-sensitive gradient coefficients. Additionally, we investigate the impact of iterative RL; for this, our experiments involve two iterations. Figure \ref{fig:iter-rl} demonstrates that iterative RL yields a marked performance gain, with the most pronounced improvement occurring after the initial iteration.

\subsubsection{Why RL Works?} In this study, we utilize reinforcement learning on a specific subset of instruction tuning data and observe a notable improvement in performance relative to the instruction-tuned baseline. To clarify the mechanism behind the effectiveness of reinforcement learning, we compare the Pass@K and Maj@K accuracy metrics for both the Instruct and RL models across two established benchmarks. As illustrated in Figure \ref{fig:maj-pass}, RL leads to higher Maj@K scores without affecting Pass@K. These results suggest that reinforcement learning primarily strengthens the model's output distribution by making it more robust, specifically by \textbf{increasing the likelihood of producing at least one correct answer in the top K outputs, rather than advancing the model’s underlying abilities.} In line with this, \citep{wang2023making} discussed a \textbf{misalignment problem} that arises in reasoning tasks for SFT models, highlighting that strategies for preference alignment can elevate the reasoning performance of SFT models \citep{yuan2023rrhf,song2023preference,wang2023making}.

\subsubsection{How to Achieve More Effective RL?}
\label{sec:effec_rl}
Our experiments show that RL can be highly effective for tasks involving mathematical reasoning. Additionally, we introduce a cohesive framework for interpreting various widely-used training strategies, framing them as either explicit or implicit instances of RL methodologies. As outlined in Equation \ref{eq:objective}, this framework is built upon three main pillars: Data Source, Algorithm, and Reward Function. We also outline several prospective research avenues related to these three foundational components.

\paragraph{Data Source} The data source constitutes the foundational input for all training paradigms. Within RL, we define the data source as comprising unlabeled questions paired with responses generated by the policy model. In this study, we restrict ourselves to utilizing questions curated during the instruction tuning phase, and generate corresponding outputs using a basic nucleus sampling approach. We hypothesize that this constraint contributes to our RL pipeline yielding improvements solely in Maj@K metrics. For subsequent research, we aim to apply our RL framework to question prompts that fall outside the original distribution, and to integrate \textbf{advanced sampling (decoding) strategies}, such as tree-search techniques \citep{yao2023tree}. Furthermore, \textbf{efficient inference techniques} \citep{xia-etal-2023-speculative,leviathan2023fast,kwon2023efficient,xia2024unlocking}, which govern how effectively the policy model explores, are also recognized as being critically important.

\paragraph{Algorithms} Algorithms utilize both data and reward signals to compute gradient coefficients, which subsequently adjust model parameters. Referencing Equation \ref{eq:objective}, the prevailing approaches largely \textbf{TRUST} the reward function, modulating the conditional probability of certain tokens in response to the reward's guidance. Nevertheless, the reward signal cannot be assumed reliable at all times, particularly when tackling highly intricate tasks. For instance, the PRM800K dataset \citep{lightman2023let}, despite being curated by expert annotators, still suffers from roughly 20\% mislabeled instances\footnote{\url{https://github.com/openai/prm800k/issues/12\#issuecomment-1728491852}}. Consequently, our focus shifts to reinforcement learning algorithms engineered for resilience in the face of noisy reward feedback. We anticipate that \textbf{WEAK-TO-STRONG} \citep{burns2023weak} alignment strategies have the potential to fundamentally reshape the landscape of learning algorithms.

\paragraph{Reward Function} The reward function serves as the principal source of training signals. Within RL, this typically takes the form of a neural reward model. We identify three pivotal avenues of research pertaining to reward models: 1) \textbf{Improving the generalization capacity of reward models.} For RL to enhance the intrinsic capabilities of LLMs—beyond merely stabilizing their existing output distributions—the reward model needs robust generalization, especially when faced with out-of-distribution prompts and complex generated responses; 2) \textbf{Capturing and representing the uncertainty inherent in reward models.} Accounting for uncertainty may play a vital role in bridging the gap between imperfect reward models and methodologies designed to leverage progressively stronger learning paradigms; 3) \textbf{Constructing high-quality process reward models efficiently} that are capable of delivering detailed, fine-grained feedback throughout the reasoning steps involved in RL training \citep{lightman2023let,wang2023math}.

\section{Conclusion, Limitation, and Future Work}  
In this work, we introduce DeepSeekMath, a model that surpasses all existing open-source systems on the competition-level MATH benchmark, achieving performance levels comparable to those of proprietary models. DeepSeekMath~is constructed by initializing with DeepSeek-Coder-v1.5 7B and is further refined through continued training over 500B tokens, notably incorporating 120B math-related tokens harvested from Common Crawl as a major component of the training corpus. Our comprehensive ablation analysis reveals that web sources present considerable opportunities for assembling high-quality mathematical datasets, whereas data from arXiv may not deliver the anticipated benefits. Furthermore, we propose Group Relative Policy Optimization (GRPO), which builds upon Proximal Policy Optimization (PPO) and offers significant improvements in mathematical reasoning abilities with reduced memory requirements. Experimental findings demonstrate the effectiveness of GRPO, even when DeepSeekMath-Instruct 7B already attains strong benchmark results. Additionally, we propose a cohesive framework that contextualizes a range of related methods, and we identify multiple promising avenues for advancing reinforcement learning techniques in this domain.

While DeepSeekMath~demonstrates strong performance on quantitative reasoning tasks, its proficiency in geometry and theorem-proving lags behind that of proprietary models. For example, during our preliminary evaluations, the model struggled with problems involving triangles and ellipses, suggesting that pre-training and fine-tuning data may exhibit selection bias. Furthermore, the model’s limited scale constrains its few-shot learning abilities; unlike GPT-4, which benefits from few-shot prompts to enhance its results, DeepSeekMath~exhibits comparable scores under both zero-shot and few-shot conditions. Moving forward, we intend to refine our data selection process to curate a higher-quality pre-training dataset. We also plan to investigate alternative approaches for reinforcement learning to bolster LLM capabilities, as discussed in Section \ref{sec:effec_rl}.

\end{document}